\par\noindent\textbf{(1)} 临界拐点条件 (∂P/∂v)\_T=(∂²P/∂v²)\_T=0。分别解 −RT/(v−b)²+2a/v³=0、2RT/(v−b)³−6a/v⁴=0，得 v\_c=3b、T\_c=8a/(27Rb)、P\_c=a/(27b²)。令 p=P/P\_c、t=T/T\_c、ν=v/v\_c，得 p=8t/(3ν−1)−3/ν²。
\par\noindent\textbf{(2)} 固定T有 dF=−Pdv，故两相Helmholtz自由能之差 F(v\_g)−F(v\_l)=−∫\_\{v\_l\}\textasciicircum{}\{v\_g\}P(v)dv。共同切线构造要求共同切线斜率为−P₀，因而 ∫\_\{v\_l\}\textasciicircum{}\{v\_g\}[P(v,T)−P₀]dv=0。这与两相压力相等及化学势相等同时成立，是Maxwell等面积而非凭图形任意选择。
\par\noindent\textbf{(3)} 共存端点由 P(v\_l,T)=P(v\_g,T)=P₀ 与积分式共同确定。对当前方程可显式积分：RT ln[(v\_g−b)/(v\_l−b)]+a(1/v\_g−1/v\_l)−P₀(v\_g−v\_l)=0。三个方程确定 v\_l,v\_g,P₀；若误取不稳定的∂P/∂v>0分支，将违反正等温压缩率。
\par\noindent\textbf{(4)} 在 t↑1 附近，以 ν=1+ω 展开约化方程，含线性温差项和 ω³非线性项；共同切线与等面积解给 ω\_g−ω\_l∝(1−t)\textasciicircum{}\{1/2\}，即平均场指数β=1/2。极近真实临界点长程涨落可改变指数，此结论是van der Waals平均场理论的内在预测。
